\chapter{ANALISIS KEAMANAN}
% Generated by scripts/render_brainrot_results.py from recorded v2 JSON.
\newcommand{\VTwoBlockRows}{%
1 bit plainteks & 256 & 49.643 & 47 & 79\\
1 bit kunci & 256 & 50.177 & 49 & 82\\
}

% Generated by scripts/render_brainrot_results.py from recorded v2 JSON.
\newcommand{\VTwoModeRows}{%
ECB & 416 & 1.931 & 50.001\\
CBC & 416 & 25.969 & 50.056\\
CFB & 404 & 24.204 & 50.000\\
OFB & 404 & 0.031 & 49.923\\
CTR & 404 & 0.031 & 49.959\\
}

% Generated by scripts/render_brainrot_results.py from recorded v2 JSON.
\newcommand{\VTwoEntropyRows}{%
ECB & 416 & 4.860 & 7.217\\
CBC & 416 & 4.860 & 7.511\\
CFB & 404 & 4.860 & 7.456\\
OFB & 404 & 4.860 & 7.463\\
CTR & 404 & 4.860 & 7.513\\
}

% Generated by scripts/render_brainrot_results.py from recorded v2 JSON.
\newcommand{\VTwoBinaryEntropyRows}{%
ECB & 1040 & 8.000 & 7.197\\
CBC & 1040 & 8.000 & 7.796\\
CFB & 1027 & 8.000 & 7.777\\
OFB & 1027 & 8.000 & 7.785\\
CTR & 1027 & 8.000 & 7.807\\
}


\section{Konfigurasi dan Reproduksibilitas}
CLI \texttt{analyze} menggunakan Brainrot v2 secara default. Masukan bawaan
adalah \path{src/testdata/fixtures/one-block.txt} sepanjang 16 byte dan
\path{src/testdata/fixtures/repeated-text.txt} sepanjang 404 byte.
Parameter bawaan adalah 256 percobaan, seed 42, serta kunci
\texttt{00112233445566778899aabbccddeeff}. Seed hanya digunakan untuk
eksperimen; IV pada CLI enkripsi dibangkitkan secara terpisah.

\begin{lstlisting}[language={},numbers=none,caption={Menjalankan analisis dari root repositori.}]
go -C src run ./cmd/analyze -cipher brainrot -modes all \
  -trials 256 -seed 42 -output ../reports/brainrot-v2/text-seed42.json
go -C src test ./...
\end{lstlisting}

Laporan JSON menyimpan ukuran blok dan kunci, seed, jumlah percobaan,
avalanche satu blok, hasil round-trip per mode, entropi, histogram 256 bin,
dan avalanche pesan utuh. Pengukuran mode hanya memakai badan cipherteks;
IV tidak masuk histogram atau penyebut persentase perubahan bit.

Hasil yang ditampilkan berasal dari \path{reports/brainrot-v2/text-seed42.json}
dan \path{reports/brainrot-v2/binary-seed42.json}. Fixture biner
\path{src/testdata/fixtures/byte-pattern.bin} terdiri atas nilai byte 0--255
yang diulang empat kali, diikuti $\mathtt{00,ff,80}$; totalnya 1027 byte.
Manifest pada folder hasil mencatat kunci, seed, ukuran dan SHA-256 fixture,
versi Go, sistem operasi, commit dasar, serta hash sumber v2. Dua eksperimen
tambahan memakai seed 1337 dengan kunci nol pada teks dan kunci semua satu
pada biner. Seluruh angka adalah pengamatan sampel, bukan ambang kelulusan
keamanan.

\section{Avalanche Effect Satu Blok}
Pada setiap percobaan, satu bit plainteks diubah dengan kunci tetap;
secara terpisah, satu bit kunci diubah dengan plainteks tetap. Percobaan
pertama memakai fixture satu blok, sedangkan berikutnya memakai blok
pseudorandom dari seed yang tercatat. Posisi bit digilir melalui seluruh
128 bit masukan; setidaknya 128 percobaan diperlukan untuk mencakup
setiap posisi. Keluaran mencatat rata-rata, minimum, maksimum,
serta rata-rata dan jumlah percobaan per posisi bit.

\begin{table}[htbp]
  \centering\small
  \caption{Hasil avalanche Brainrot v2 satu blok, 256 percobaan dan seed 42.}
  \begin{tabularx}{\linewidth}{@{}lYYYY@{}}
    \toprule Perubahan & Percobaan & Rata-rata (\%) & Minimum (bit) & Maksimum (bit)\\
    \midrule
    \VTwoBlockRows
    \bottomrule
  \end{tabularx}
\end{table}

Dengan 256 percobaan, setiap posisi bit diuji dua kali. Rata-rata perubahan
bit plainteks dan kunci mendekati separuh keluaran pada sampel ini.
\Isi{Kembangkan pembahasan variasi per posisi bit dari JSON dan
bandingkan dengan eksperimen kunci/seed lainnya.}

\section{Avalanche Effect pada Seluruh Mode}
Untuk pesan utuh, posisi bit plainteks dipilih pseudorandom dan posisi
bit kunci digilir. Pasangan enkripsi sebelum dan sesudah perubahan
menggunakan IV yang sama agar perbedaan berasal dari bit yang diubah.
Penyebutnya adalah jumlah bit seluruh badan cipherteks, termasuk padding
ECB/CBC, sehingga berbeda dari pengujian satu blok.

\begin{table}[htbp]
  \centering\small
  \caption{Avalanche pesan teks 404 byte, 256 percobaan dan seed 42.}
  \begin{tabularx}{\linewidth}{@{}lYYY@{}}
    \toprule Mode & Panjang badan (byte) & Perubahan plainteks (\%) & Perubahan kunci (\%)\\
    \midrule
    \VTwoModeRows
    \bottomrule
  \end{tabularx}
\end{table}
Pada OFB dan CTR dengan kunci serta IV tetap, satu bit plainteks yang
diubah menghasilkan satu bit cipherteks yang berubah karena operasi XOR
dengan keystream yang sama. Pada ECB perubahan terbatas pada blok terkait;
pada CBC dan CFB pengaruhnya bergantung pada posisi perubahan dan umpan
balik. Interpretasi hasil harus membedakan sifat mode dari difusi cipher
satu blok.

\Isi{Bahas hasil kelima mode, posisi bit yang dipilih, pengaruh panjang
pesan, dan perbedaan perubahan plainteks terhadap perubahan kunci.}

\section{Analisis Entropi}
\begin{table}[htbp]
  \centering\small
  \caption{Entropi empiris untuk fixture teks, seed 42; ukuran adalah badan cipherteks.}
  \begin{tabularx}{\linewidth}{@{}lYYY@{}}
    \toprule Mode & Ukuran sampel (byte) & Plainteks (bit/byte) & Cipherteks (bit/byte)\\
    \midrule
    \VTwoEntropyRows
    \bottomrule
  \end{tabularx}
\end{table}
\begin{table}[htbp]
  \centering\small
  \caption{Entropi empiris untuk fixture biner, seed 42.}
  \begin{tabularx}{\linewidth}{@{}lYYY@{}}
    \toprule Mode & Ukuran badan (byte) & Plainteks (bit/byte) & Cipherteks (bit/byte)\\
    \midrule
    \VTwoBinaryEntropyRows
    \bottomrule
  \end{tabularx}
\end{table}
\Isi{Jelaskan pengaruh ukuran sampel, pola plainteks, dan padding.
Nilai mendekati 8 bit/byte saja tidak membuktikan keamanan kriptografis.}

\section{Analisis Histogram}
\begin{figure}[htbp]
  \centering
  \GambarAtauPlaceholder{histogram-semua-mode.png}{\linewidth}{11cm}
  \caption{Histogram fixture teks dan cipherteks lima mode pada seed 42;
    sumbu y memakai skala berbeda dan garis putus-putus menandai distribusi seragam.}
  \label{fig:histogram}
\end{figure}
\begin{figure}[htbp]
  \centering
  \GambarAtauPlaceholder{histogram-biner.png}{\linewidth}{11cm}
  \caption{Histogram fixture biner dan cipherteks lima mode, seed 42;
    sumbu y memakai skala berbeda.}
\end{figure}
Grafik memakai proporsi byte agar panjang badan yang berbeda dapat
dibandingkan. Skala y berbeda antar panel dan tercantum pada setiap sumbu.
\Isi{Bahas pola berulang, khususnya ECB, serta kedekatan distribusi empiris
terhadap distribusi seragam dengan memperhatikan ukuran sampel.}

\section{Pengujian Tambahan dan Keterbatasan}
\Isi{Jika dilakukan, tambahkan pengujian ukuran kecil/sedang/besar,
waktu proses, penggunaan memori, perubahan per putaran melalui Tracer,
atau kajian struktur S-box dan matriks. Sertakan metode serta hasil aktual.}

Pengujian otomatis mempertahankan \texttt{TEST-ONLY-XOR} untuk hasil
chaining yang dapat diprediksi, dan menambahkan pengujian khusus Brainrot
v2. Empat vektor referensi independen memeriksa cipherteks, 21 subkunci
per kasus, serta 66 snapshot. Invers, input berpola dan pseudorandom,
overlap buffer, pemetaan semua 128 bit BPT, serta seluruh 65.536 pasangan
byte/subkunci untuk penjumlahan diperiksa. Lima mode juga diuji dengan
Brainrot pada masukan kosong, batas blok, padding tidak valid, dan overflow CTR.

Pemeriksaan rank biner menguji 36.416 pasangan support untuk setiap
jenis propagasi dan membuktikan branch number byte minimal 5 pada
lapisan ShiftRows67--MixColumns67--BPT revisi untuk perbedaan dan masker
linier. Referensi Python memeriksa batas yang sama secara independen.
Batas tersebut tidak mencakup FCM, penjumlahan, atau seluruh cipher.

Benchmark tiga kali pada Linux amd64, AMD Ryzen 5 5600H, menghasilkan
median enkripsi 20.225 ns/blok pada cipher lama dan 24.959 ns/blok pada
v2 (kenaikan waktu sekitar 23,4\%). Median dekripsi berubah dari
20.141 menjadi 24.673 ns/blok (sekitar 22,5\%). Tidak ada alokasi heap
per blok pada kedua versi. Angka tersebut berasal dari benchmark lokal
dan bukan batas performa universal; keluaran mentah disimpan bersama hasil.

FCM memiliki simetri rotasi pada kedua half, dan fungsi internalnya
memiliki empat kemungkinan perbedaan keluaran dengan probabilitas
$1/4$ untuk perbedaan masukan satu bit pada input seragam. Penjumlahan
byte mempertahankan perbedaan XOR bit tertinggi untuk setiap subkunci.
Pasangan subkunci penjumlahan/XOR $(a,b)$ dan
$((a+128)\bmod256,b\oplus128)$ juga menghasilkan pemetaan byte sama.
Ini adalah sifat komponen, bukan bukti serangan pada cipher lengkap
atau adanya kunci utama ekuivalen.

\Isi{Nyatakan cakupan pengujian akhir, keterbatasan sampel, dan analisis
yang belum dilakukan. Kaitkan kesimpulan hanya dengan bukti yang tersedia.}
